\section{Première semaine : Conceptualisation de l'atelier}
  Durant les premiers jours, j'ai été accueilli par l'équipe pédagogique de l'\enseirb et j'ai eu l'occasion de rencontrer les autres étudiants de première année qui participaient également à l'atelier. Nous avons été introduits aux objectifs de l'atelier et sur les attentes du stage des élèves de seconde. \par
 Afin de préparer l'atelier, nous avons régulièrement eu des réunions avec notre tuteur de stage Charles Brazier, chercheur au LaBri et professeur à l'\enseirb. Ces réunions ont été l'occasion de discuter des différentes idées et approches pour l'atelier, ainsi que de recevoir des conseils sur la manière d'expliquer le fonctionnement d'un modèle de langage\footnote{Modèle d'intelligence artificielle s'exprimant dans un langage textuel.}. \par
 Plusieurs idées ont été proposées pour l'atelier, mais la plupart nous ont paru assez peu pertinentes ou abordables pour des élèves de seconde. En effet, il a été difficile de trouver un sujet qui soit à la fois intéressant et compréhensible pour des élèves de cet âge. La connaissance que j'ai acquise grâce à mon expérience de soutien scolaire m'a permis de savoir précisément ce que les élèves pouvaient aborder à leur niveau. Après plusieurs discussions, nous avons finalement décidé de nous concentrer sur un atelier axé en deux parties~:
\begin{itemize}
  \item Une première partie consiste à faire réfléchir les élèves sur le fonctionnement primitif d'un modèle de langage, en transposant les problèmes rencontrés à des problèmes du quotidien comme "Deviner la suite d'une suite de symboles".

  \item La seconde partie consiste à introduire aux élèves à l'utilisation d'un modèle de langage, d'abord simplement en leur laissant parler avec celui-ci, puis en leur faisant découvrir les limites de ce modèle et les biais qu'il peut contenir. 
\end{itemize}

